\documentclass[12pt]{amsart}
% Default margins are too wide all the way around. I reset them here
\setlength{\hoffset}{-.75in}
\setlength{\textwidth}{6.5in}
\setlength{\voffset}{-.5in}
\setlength{\textheight}{9.0in}
\usepackage{amsmath}
\usepackage{amssymb}
\usepackage{amsthm}
\usepackage{pgfplots}
\usepackage{enumerate}
\usepackage{tikz}
\usetikzlibrary{shadows,arrows,arrows.meta,shapes,positioning,calc}% Required for the Circle and Triangle arrow tips
\usepackage{setspace}
\usepackage{siunitx}
\usepackage{enumitem}
\usepackage{tcolorbox}
\usepackage{array}
\usepackage{microtype}

\theoremstyle{conjecture}
\newtheorem{conjecture}{Conjecture}

\theoremstyle{definition}
\newtheorem{definition}{Definition}

\theoremstyle{question}
\newtheorem{question}{Question}

\theoremstyle{theorem}
\newtheorem{theorem}{Theorem}

\newenvironment{solution}
{\begin{proof}[Solution]}
{\end{proof}}

\newcommand{\sectionrule}{%
    \vspace{0.3em}
    \hrule
    \vspace{0.8em}
}

\title{Dynamical Systems:\\Class August 27, 2026\\ 3.a Nondimensionalization.}
\author{Rob Tompkins}
\date{\today}


\begin{document}

    \maketitle

    \section*{Nondimensionalization}

    \subsection*{Problem 1}

    Let
    \[
        \dot N=rN\left(1-\frac{N}{K}\right)-H.
    \]
    This is a logistic population model where there is a constant
    harvesting rate reducing population growth.

    \begin{enumerate}[label=(\alph*)]

% ------------------------------------------------------------------
        \item
        For each of the variables and each of the parameters, identify its
        associated dimension. Write this out in expressions of the form
        \([a]=L\).

        \subsubsection*{Solution}

        Let \(P\) denote the dimension of population and let \(T\) denote the
        dimension of time. Since \(N\) is a population and \(t\) is time,
        \[
            [N]=P,
            \qquad
            [t]=T.
        \]

        The left-hand side has dimension
        \[
            [\dot N]
            =
            \left[\frac{dN}{dt}\right]
            =
            \frac{P}{T}.
        \]

        The logistic term must therefore also have dimension \(P/T\):
        \[
            \left[
                rN\left(1-\frac{N}{K}\right)
            \right]
            =
            \frac{P}{T}.
        \]

        The quantity inside the parentheses must be dimensionless. Therefore,
        \[
            \left[\frac{N}{K}\right]=1,
        \]
        so
        \[
            [K]=P.
        \]

        It follows that
        \[
            [r][N]=\frac{P}{T},
        \]
        and hence
        \[
            [r]=T^{-1}.
        \]

        Finally, the harvesting term \(H\) must have the same dimension as
        \(\dot N\), so
        \[
            [H]=\frac{P}{T}.
        \]

        Thus,
        \[
            \boxed{
                [N]=P,\qquad
                [t]=T,\qquad
                [r]=T^{-1},\qquad
                [K]=P,\qquad
                [H]=PT^{-1}.
            }
        \]

% ------------------------------------------------------------------
        \item
        Create dimensional constants \(N_0\) and \(T_0\), and use them to
        create nondimensional variables
        \[
            x=\frac{N}{N_0}
            \qquad\text{and}\qquad
            \tau=\frac{t}{T_0}.
        \]
        Substitute \(x\) and \(\tau\) into the equation and simplify.

        \subsubsection*{Solution}

        Define
        \[
            x=\frac{N}{N_0},
            \qquad
            \tau=\frac{t}{T_0}.
        \]

        Thus,
        \[
            N=N_0x,
            \qquad
            t=T_0\tau.
        \]

        We transform the derivative using the chain rule:
        \[
            \frac{dN}{dt}
            =
            \frac{d(N_0x)}{d\tau}
            \frac{d\tau}{dt}.
        \]

        Since
        \[
            \frac{d\tau}{dt}=\frac{1}{T_0},
        \]
        we obtain
        \[
            \frac{dN}{dt}
            =
            \frac{N_0}{T_0}\frac{dx}{d\tau}.
        \]

        Substituting \(N=N_0x\) into the original equation gives
        \[
            \frac{N_0}{T_0}\frac{dx}{d\tau}
            =
            rN_0x
            \left(
                1-\frac{N_0x}{K}
            \right)
            -H.
        \]

        Multiplying through by \(T_0/N_0\),
        \[
            \frac{dx}{d\tau}
            =
            rT_0x
            \left(
                1-\frac{N_0}{K}x
            \right)
            -
            \frac{HT_0}{N_0}.
        \]

        Therefore,
        \[
            \boxed{
                \frac{dx}{d\tau}
                =
                rT_0x
                \left(
                    1-\frac{N_0}{K}x
                \right)
                -
                \frac{HT_0}{N_0}.
            }
        \]

        Equivalently, expanding the logistic term,
        \[
            \boxed{
                \frac{dx}{d\tau}
                =
                rT_0x
                -
                \frac{rT_0N_0}{K}x^2
                -
                \frac{HT_0}{N_0}.
            }
        \]

% ------------------------------------------------------------------
        \item
        List all nondimensional groups that arise. Consider a combination a
        nondimensional group if the combination is nondimensional but every
        piece would be dimensional if it were broken apart in any way.

        \subsubsection*{Solution}

        From
        \[
            \frac{dx}{d\tau}
            =
            rT_0x
            \left(
                1-\frac{N_0}{K}x
            \right)
            -
            \frac{HT_0}{N_0},
        \]
        the nondimensional groups are
        \[
            rT_0,
            \qquad
            \frac{N_0}{K},
            \qquad
            \frac{HT_0}{N_0}.
        \]

        We may verify each directly.

        First,
        \[
            [rT_0]
            =
            T^{-1}T
            =
            1.
        \]

        Second,
        \[
            \left[\frac{N_0}{K}\right]
            =
            \frac{P}{P}
            =
            1.
        \]

        Third,
        \[
            \left[\frac{HT_0}{N_0}\right]
            =
            \frac{(P/T)T}{P}
            =
            1.
        \]

        Thus,
        \[
            \boxed{
                rT_0,\qquad
                \frac{N_0}{K},\qquad
                \frac{HT_0}{N_0}
            }
        \]
        are the nondimensional groups that arise.

% ------------------------------------------------------------------
        \item
        Make choices for values of the constants \(T_0\) and \(N_0\) that
        eliminate two of the nondimensional groups.

        \subsubsection*{Solution}

        A natural choice is to make
        \[
            rT_0=1.
        \]

        Therefore choose
        \[
            \boxed{T_0=\frac{1}{r}}.
        \]

        We can also choose
        \[
            \frac{N_0}{K}=1,
        \]
        which gives
        \[
            \boxed{N_0=K}.
        \]

        With these choices,
        \[
            rT_0=1
            \qquad\text{and}\qquad
            \frac{N_0}{K}=1.
        \]

        The differential equation becomes
        \[
            \frac{dx}{d\tau}
            =
            x(1-x)
            -
            \frac{H}{rK}.
        \]

% ------------------------------------------------------------------
        \item
        Define a new nondimensional parameter (use a Greek letter such as
        \(\alpha,\beta,\gamma,\mu\)), and rewrite your equation as a
        nondimensional one.

        \subsubsection*{Solution}

        Define
        \[
            \boxed{
                \alpha=\frac{H}{rK}
            }.
        \]

        The nondimensional equation becomes
        \[
            \boxed{
                \frac{dx}{d\tau}
                =
                x(1-x)-\alpha.
            }
        \]

        Here
        \[
            x=\frac{N}{K},
            \qquad
            \tau=rt,
            \qquad
            \alpha=\frac{H}{rK}.
        \]

% ------------------------------------------------------------------
        \item
        How many parameters are there in the nondimensional system? How does
        this compare to the number in the dimensional version? Notice that
        each dimensional constant we introduce enables us to remove a
        parameter, so that the nondimensional equation has fewer parameters
        than the dimensional one. This result on the reduction in the number
        of parameters from nondimensionalization is called the Buckingham Pi
        theorem. It is called the ``pi'' theorem because he used the symbols
        \(\Pi_1,\Pi_2,\ldots\) to represent the nondimensional groups.

        \subsubsection*{Solution}

        The original dimensional equation
        \[
            \dot N
            =
            rN\left(1-\frac{N}{K}\right)-H
        \]
        contains three parameters:
        \[
            r,\qquad K,\qquad H.
        \]

        After nondimensionalization, the equation is
        \[
            \frac{dx}{d\tau}
            =
            x(1-x)-\alpha,
        \]
        which contains only one parameter:
        \[
            \alpha=\frac{H}{rK}.
        \]

        Therefore,
        \[
            \boxed{
                \text{dimensional model: 3 parameters}
            }
        \]
        while
        \[
            \boxed{
                \text{nondimensional model: 1 parameter}.
            }
        \]

        We introduced two dimensional scaling constants, \(T_0\) and \(N_0\),
        and these allowed us to eliminate two of the original independent
        parameters.

        Thus the parameter count has been reduced from
        \[
            \boxed{3\longrightarrow 1}.
        \]

    \end{enumerate}


    \newpage

% ==================================================================
% PROBLEM 2
% ==================================================================

    \subsection*{Problem 2}

    \textit{Rotate your roles: questioner \(\rightarrow\) timekeeper
        \(\rightarrow\) scribe.}

    Let
    \[
        \dot N
        =
        rN\left(1-\frac{N}{K}\right)
        -
        \frac{HN}{A+N}.
    \]
    This is a slightly different harvesting case.

    \begin{enumerate}[label=(\alph*)]

% ------------------------------------------------------------------
        \item
        This harvesting term,
        \[
            \frac{HN}{A+N},
        \]
        is in the form of a special function called a Monod function. Plot an
        approximation of the Monod function by hand without using any plotting
        tools.

        \subsubsection*{Solution}

        Consider the harvesting function
        \[
            F(N)=\frac{HN}{A+N}.
        \]

        At \(N=0\),
        \[
            F(0)=0.
        \]

        At \(N=A\),
        \[
            F(A)
            =
            \frac{HA}{A+A}
            =
            \frac{H}{2}.
        \]

        As \(N\to\infty\),
        \[
            \lim_{N\to\infty}
            \frac{HN}{A+N}
            =
            H.
        \]

        Thus the horizontal asymptote is
        \[
            \boxed{F(N)=H}.
        \]

        For small \(N\), we have
        \[
            A+N\approx A,
        \]
        so
        \[
            F(N)
            \approx
            \frac{H}{A}N.
        \]

        Therefore, near \(N=0\), the function is approximately linear with
        slope
        \[
            \frac{H}{A}.
        \]

        A sketch of the function is:

        \begin{center}
            \begin{tikzpicture}
                \begin{axis}[
                    width=11cm,
                    height=6.5cm,
                    axis lines=middle,
                    xlabel={$N$},
                ylabel={$F(N)$},
                xmin=0,
                xmax=5,
                ymin=0,
                ymax=1.2,
                xtick={1},
                xticklabels={$A$},
                ytick={0.5,1},
                yticklabels={$\frac{H}{2}$,$H$},
                samples=200,
                domain=0:5,
                clip=false
                ]
                    \addplot[thick] {x/(1+x)};

                    \addplot[dashed,domain=0:5] {1};

                    \addplot[only marks,mark=*]
                    coordinates {(1,0.5)};

                    \node at (axis cs:1.2,0.43)
                        {$\left(A,\frac{H}{2}\right)$};

                \end{axis}
            \end{tikzpicture}
        \end{center}

        The important qualitative features are therefore
        \[
            \boxed{
                F(0)=0,\qquad
                F(A)=\frac{H}{2},\qquad
                F(N)\to H
                \text{ as }N\to\infty.
            }
        \]

% ------------------------------------------------------------------
        \item
        What do you think of this function as a description of a harvesting
        process?

        \subsubsection*{Solution}

        The function
        \[
            F(N)=\frac{HN}{A+N}
        \]
        gives a harvesting rate that depends on the available population.

        When the population is small,
        \[
            F(N)\approx\frac{H}{A}N,
        \]
        so the harvesting rate is also small. Thus, when there are few
        individuals available, relatively few individuals are harvested.

        As the population increases, the harvesting rate increases, but it
        does not increase without bound. Instead,
        \[
            F(N)\to H.
        \]

        Thus \(H\) represents the maximum harvesting rate.

        This is a reasonable model when the ability to harvest is limited by
        the availability of the population at low population levels and by
        some maximum harvesting capacity at high population levels.

% ------------------------------------------------------------------
        \item
        Identify the dimension of each of the variables and parameters. Once
        you nondimensionalize how many parameters do you expect to remain?

        \subsubsection*{Solution}

        Again let \(P\) denote population and \(T\) denote time.

        We have
        \[
            [N]=P,
            \qquad
            [t]=T.
        \]

        As before,
        \[
            [r]=T^{-1},
            \qquad
            [K]=P.
        \]

        Because \(A\) is added to \(N\),
        \[
            [A]=[N]=P.
        \]

        Now consider
        \[
            \frac{HN}{A+N}.
        \]

        Since
        \[
            \left[
                \frac{N}{A+N}
            \right]
            =1,
        \]
        the parameter \(H\) itself must have the dimension of population per
        unit time:
        \[
            [H]=\frac{P}{T}.
        \]

        Therefore,
        \[
            \boxed{
                [N]=P,\qquad
                [t]=T,\qquad
                [r]=T^{-1},\qquad
                [K]=P,\qquad
                [A]=P,\qquad
                [H]=PT^{-1}.
            }
        \]

        There are four dimensional parameters:
        \[
            r,\qquad K,\qquad A,\qquad H.
        \]

        There are two fundamental dimensions involved:
        \[
            P
            \qquad\text{and}\qquad
            T.
        \]

        We therefore expect nondimensionalization to reduce the number of
        independent parameters by two. Hence we expect
        \[
            \boxed{4-2=2}
        \]
        independent nondimensional parameters to remain.

% ------------------------------------------------------------------
        \item
        Nondimensionalize this equation.

        As part of this process, identify the dimensionless groups and check
        your work with another team.

        There are multiple good choices for \(N_0\). What are some reasons to
        choose one or the other?

        \subsubsection*{Solution}

        Introduce
        \[
            x=\frac{N}{N_0},
            \qquad
            \tau=\frac{t}{T_0}.
        \]

        Thus
        \[
            N=N_0x,
            \qquad
            t=T_0\tau.
        \]

        As before,
        \[
            \frac{dN}{dt}
            =
            \frac{N_0}{T_0}\frac{dx}{d\tau}.
        \]

        Substituting into
        \[
            \dot N
            =
            rN\left(1-\frac{N}{K}\right)
            -
            \frac{HN}{A+N}
        \]
        gives
        \[
            \frac{N_0}{T_0}\frac{dx}{d\tau}
            =
            rN_0x
            \left(
                1-\frac{N_0x}{K}
            \right)
            -
            \frac{HN_0x}{A+N_0x}.
        \]

        Multiplying through by \(T_0/N_0\),
        \[
            \frac{dx}{d\tau}
            =
            rT_0x
            \left(
                1-\frac{N_0}{K}x
            \right)
            -
            \frac{HT_0x}{A+N_0x}.
        \]

        Factor \(N_0\) from the denominator of the harvesting term:
        \[
            A+N_0x
            =
            N_0
            \left(
                \frac{A}{N_0}+x
            \right).
        \]

        Therefore,
        \[
            \boxed{
                \frac{dx}{d\tau}
                =
                rT_0x
                \left(
                    1-\frac{N_0}{K}x
                \right)
                -
                \frac{HT_0}{N_0}
                \frac{x}{A/N_0+x}.
            }
        \]

        The dimensionless groups that arise are
        \[
            \boxed{
                rT_0,\qquad
                \frac{N_0}{K},\qquad
                \frac{A}{N_0},\qquad
                \frac{HT_0}{N_0}.
            }
        \]

        These groups are not all independent because the freely chosen scales
        \(N_0\) and \(T_0\) can be used to eliminate two of them.

        A convenient time scale is
        \[
            \boxed{
                T_0=\frac{1}{r}.
            }
        \]

        For the population scale there are at least two especially natural
        choices.

        \paragraph{Choice 1: \(N_0=K\).}

        If
        \[
            N_0=K,
        \]
        then \(x=N/K\) measures population relative to the carrying capacity.
        The equation becomes
        \[
            \frac{dx}{d\tau}
            =
            x(1-x)
            -
            \frac{H}{rK}
            \frac{x}{A/K+x}.
        \]

        Define
        \[
            \alpha=\frac{H}{rK},
            \qquad
            \beta=\frac{A}{K}.
        \]

        Then
        \[
            \boxed{
                \frac{dx}{d\tau}
                =
                x(1-x)
                -
                \alpha\frac{x}{\beta+x}.
            }
        \]

        This choice is useful because the carrying capacity occurs at
        \[
            x=1.
        \]

        \paragraph{Choice 2: \(N_0=A\).}

        Alternatively, let
        \[
            \boxed{N_0=A}.
        \]

        This is especially convenient for studying the Monod harvesting term.
        The equation becomes
        \[
            \frac{dx}{d\tau}
            =
            x
            \left(
                1-\frac{A}{K}x
            \right)
            -
            \frac{H}{rA}
            \frac{x}{1+x}.
        \]

        Define
        \[
            \alpha=\frac{A}{K},
            \qquad
            \beta=\frac{H}{rA}.
        \]

        Then
        \[
            \boxed{
                \frac{dx}{d\tau}
                =
                x(1-\alpha x)
                -
                \beta\frac{x}{1+x}.
            }
        \]

        This form contains only two nondimensional parameters,
        \[
            \boxed{
                \alpha=\frac{A}{K},
                \qquad
                \beta=\frac{H}{rA}.
            }
        \]

        Thus, as predicted, four dimensional parameters have been reduced to
        two nondimensional parameters.

        Choosing \(N_0=K\) is convenient when the carrying capacity is the
        important population scale. Choosing \(N_0=A\) is convenient when the
        harvesting function is the primary object of interest, because the
        Monod function then takes the particularly simple form
        \[
            \frac{x}{1+x}.
        \]

% ------------------------------------------------------------------
        \item
        Now that it is nondimensional, take another look at the expression for
        your harvesting function. Plot it using appropriate axis ticks. How
        did your axis tick labels change?

        \subsubsection*{Solution}

        Using the choice
        \[
            N_0=A,
        \]
        we have
        \[
            x=\frac{N}{A}.
        \]

        The original harvesting function is
        \[
            F(N)=\frac{HN}{A+N}.
        \]

        Divide the harvesting rate by its natural scale \(H\). Define
        \[
            f(x)=\frac{F(N)}{H}.
        \]

        Since \(N=Ax\),
        \[
            f(x)
            =
            \frac{1}{H}
            \frac{H(Ax)}{A+Ax}.
        \]

        Therefore,
        \[
            f(x)
            =
            \frac{Ax}{A(1+x)}
            =
            \boxed{
                \frac{x}{1+x}
            }.
        \]

        The dimensional function
        \[
            F(N)=\frac{HN}{A+N}
        \]
        has the important point
        \[
            \left(A,\frac{H}{2}\right)
        \]
        and horizontal asymptote
        \[
            F=H.
        \]

        After nondimensionalization these become
        \[
            \boxed{
                \left(1,\frac{1}{2}\right)
            }
        \]
        and
        \[
            \boxed{
                f=1.
            }
        \]

        Thus all of the dimensional labels \(A\) and \(H\) disappear from the
        axes.

        The nondimensional plot is:

        \begin{center}
            \begin{tikzpicture}
                \begin{axis}[
                    width=11cm,
                    height=6.5cm,
                    axis lines=middle,
                    xlabel={$x=N/A$},
                ylabel={$f(x)=F/H$},
                xmin=0,
                xmax=5,
                ymin=0,
                ymax=1.2,
                xtick={0,1,2,3,4},
                ytick={0,0.5,1},
                yticklabels={$0$,$\frac12$,$1$},
                samples=200,
                domain=0:5,
                clip=false
                ]

                    \addplot[thick] {x/(1+x)};

                    \addplot[dashed,domain=0:5] {1};

                    \addplot[only marks,mark=*]
                    coordinates {(1,0.5)};

                    \node at (axis cs:1.45,0.43)
                        {$\left(1,\frac12\right)$};

                \end{axis}
            \end{tikzpicture}
        \end{center}

        Originally, increments on the horizontal axis were measured in units
        of \(A\):
        \[
            0,\ A,\ 2A,\ 3A,\ldots
        \]
        and vertical values were measured in units of \(H\).

        After nondimensionalization, the corresponding horizontal ticks are
        simply
        \[
            0,\ 1,\ 2,\ 3,\ldots
        \]
        and the vertical asymptote value changes from the dimensional value
        \(H\) to the dimensionless value
        \[
            1.
        \]

        In particular,
        \[
            \boxed{
                N=A
                \quad\longrightarrow\quad
                x=1
            }
        \]
        and
        \[
            \boxed{
                F(A)=\frac{H}{2}
                \quad\longrightarrow\quad
                f(1)=\frac12.
            }
        \]

    \end{enumerate}

\end{document}